\subsection{Learnable hybrid basis}
\label{sec:track-c}

The basis ablation (Section~\ref{sec:basis-results}) found B-spline (compact
local support, $3.9\times$ RBF's cost) and RBF (smooth global support, cheap)
statistically tied on accuracy. We therefore train the hybrid edge of
Section~\ref{sec:hybrid-method} (Equation~\ref{eq:hybrid-basis}) on
Lotka-Volterra and on the fixed pendulum (Section~\ref{sec:crossdomain}). If
one global blend systematically beat both pure bases, the basis choice in
Equation~\ref{eq:kan-edge} would be worth learning rather than fixing by
hand.

\begin{expectbox}
\expectrow{Question}{Can a network do better by learning its own blend of B-spline and RBF, rather than committing to one basis in advance?}
\expectrow{Expected}{A learned blend should match or beat the better pure basis, by adaptively favoring local support where useful and global support elsewhere.}
\expectrow{Observed}{Hybrid never beats the better pure basis on either system. On the pendulum, it actively overfits past 3000 epochs in a way pure RBF does not.}
\end{expectbox}

\begin{table*}[t]
\centering
\small
\renewcommand{\arraystretch}{1.15}
\caption{Hybrid basis vs.\ pure bases, full 10{,}000-epoch budget.}
\label{tab:hybrid-basis}
\begin{tabularx}{\textwidth}{@{}L{5.0cm} C{2.6cm} C{2.8cm} C{2.0cm} X@{}}
\toprule
\hdrrow \thh{System / basis} & \thh{Training MSE} & \thh{Extrapolation MSE} & \thh{$R^2$} & \thh{Final gate}\\
\midrule
Lotka-Volterra -- RBF (pure) & $8.83\times10^{-5}$ & $8.92\times10^{-5}$ & $1.0000$ & --\\
Lotka-Volterra -- B-spline (pure) & $7.20\times10^{-5}$ & $8.38\times10^{-5}$ & $1.0000$ & --\\
Lotka-Volterra -- Hybrid & $1.47\times10^{-4}$ & $1.78\times10^{-4}$ & $0.9999$ & $\alpha{=}0.12,\beta{=}0.88$\\
\addlinespace[2pt]
Pendulum -- RBF (pure) & $9.28\times10^{-5}$ & $0.090$ & $+0.647$ & --\\
Pendulum -- Hybrid (10,000 epochs) & $9.44\times10^{-5}$ & $\mathbf{0.809}$ & $\mathbf{-2.17}$ & $\alpha{=}0.02,\beta{=}0.98$\\
\bottomrule
\end{tabularx}
\end{table*}

\begin{figure}[tb]
\centering
\includegraphics[width=\columnwidth]{figures/track_c/lv_full_alpha_beta.png}\\[4pt]
\includegraphics[width=\columnwidth]{figures/track_c/pendulum_full_alpha_beta.png}
\caption{Learned $\alpha$/$\beta$ gate evolution: Lotka-Volterra settles into a
genuine partial blend (top); the pendulum gate reaches $\approx\!99.5\%$
RBF by epoch 250 and ends at $98\%$ (bottom); it settles well before the
overfitting onset discussed below.}
\label{fig:hybrid-gates}
\end{figure}

\paragraph{Lotka-Volterra: no benefit, real cost.} The gate settles into a
genuine partial blend, retaining $\approx\!12\%$ B-spline, which confirms the
mechanism works as designed. Hybrid ties both pure bases on final accuracy
while paying roughly the \emph{sum} of their per-epoch compute.

\paragraph{Pendulum: overfitting emerges only past a specific budget.} A
short probe run at 2000 epochs first surfaced a practical wrinkle. The gate
moved in the correct direction (toward RBF) but too slowly, so the network
still carried meaningful B-spline contamination at an epoch count where the
pure-RBF pendulum recipe was already nearly converged. Giving the gate its
own, faster learning rate ($15\times$ the base rate, direction untouched)
closed nearly all of that gap at the same probe budget
(Figure~\ref{fig:hybrid-probe}, top), and we carried this into the full run.

\begin{figure}[tb]
\centering
\includegraphics[width=\columnwidth]{figures/track_c/hybrid_probe_pendulum_alpha_beta.png}\\[4pt]
\includegraphics[width=\columnwidth]{figures/track_c/hybrid_pendulum_3k_alpha_beta.png}
\caption{Top: the 2000-epoch probe run's gate trace, showing the slow drift
toward RBF that motivated the gate-learning-rate fix. Bottom: the gate trace
for the 3000-epoch pre-overfitting-optimum run discussed below, settling to
$\beta\approx0.996$.}
\label{fig:hybrid-probe}
\end{figure}

At the full 10{,}000-epoch budget, hybrid matches pure RBF's training fit
almost exactly, but its extrapolation collapses ($R^2=-2.17$). The monitor
loss bottoms out near epoch 2950 ($\approx0.054$), then rises to a
$\approx0.4$ plateau for the remaining 7000 epochs while training loss keeps
improving. This is textbook overfitting, visible directly in the saved
training history. The gate settles well before this onset: $\beta$ reaches
$\approx0.995$ by epoch 250, some 2750 epochs before overfitting begins, and
pure RBF trained on the identical recipe shows no such pattern at all.
Capping training at the pre-overfit optimum (epoch 3000, chosen from the
10,000-epoch run's own point of minimum monitor loss) shows hybrid
\emph{ahead} of pure RBF at the same epoch count ($0.055$ vs.\ $0.081$
full-horizon MSE) with extrapolation $R^2=+0.580$, approaching pure RBF's
fully-converged number ($+0.647$). Figure~\ref{fig:hybrid-pendulum-runs} shows
both runs side by side: the loss curves of the full run make the overfitting
onset visible, and its phase portrait shows the extrapolation collapsing into
a small wrong loop, while the 3000-epoch run still follows the decaying
spiral.

\begin{figure*}[t]
\centering
\begin{minipage}[c]{0.56\textwidth}
\centering
\includegraphics[width=\linewidth]{figures/track_c/pendulum_full_loss_curves.png}
\end{minipage}\hfill
\begin{minipage}[c]{0.40\textwidth}
\centering
\includegraphics[width=\linewidth]{figures/track_c/pendulum_full_phase_space.png}
\end{minipage}\\[6pt]
\begin{minipage}[c]{0.56\textwidth}
\centering
\includegraphics[width=\linewidth]{figures/track_c/pendulum_3k_loss_curves.png}
\end{minipage}\hfill
\begin{minipage}[c]{0.40\textwidth}
\centering
\includegraphics[width=\linewidth]{figures/track_c/pendulum_3k_phase_space.png}
\end{minipage}
\caption{Hybrid basis on the pendulum. Top: the full 10{,}000-epoch
run; training loss keeps falling while the extrapolation (monitor) loss
bottoms out near epoch 2950 and climbs back to $\approx0.4$ (left), and the
extrapolated orbit collapses into a small wrong loop (right). Bottom: the same
recipe stopped at 3000 epochs; train and monitor loss fall together (left)
and the extrapolation follows the decaying spiral (right).}
\label{fig:hybrid-pendulum-runs}
\end{figure*}

\begin{keypoint}[Verdict]
The learnable hybrid basis never outperforms the better pure basis on either
system at any budget tested. Its gate mechanism works exactly as designed
(gradient-driven, converging differently and sensibly per system), but that
alone does not buy an accuracy edge. On the pendulum specifically, a
network that \emph{starts} as a blend generalizes measurably worse over long
training than one that starts pure RBF, despite an almost-identical final
gate and matching training accuracy.
\end{keypoint}
